% Options for packages loaded elsewhere
\PassOptionsToPackage{unicode}{hyperref}
\PassOptionsToPackage{hyphens}{url}
%
\documentclass[
  11pt,
]{ctexart}
\usepackage{amsmath,amssymb}
\usepackage{iftex}
\ifPDFTeX
  \usepackage[T1]{fontenc}
  \usepackage[utf8]{inputenc}
  \usepackage{textcomp} % provide euro and other symbols
\else % if luatex or xetex
  \usepackage{unicode-math} % this also loads fontspec
  \defaultfontfeatures{Scale=MatchLowercase}
  \defaultfontfeatures[\rmfamily]{Ligatures=TeX,Scale=1}
\fi
\usepackage{lmodern}
\ifPDFTeX\else
  % xetex/luatex font selection
\fi
% Use upquote if available, for straight quotes in verbatim environments
\IfFileExists{upquote.sty}{\usepackage{upquote}}{}
\IfFileExists{microtype.sty}{% use microtype if available
  \usepackage[]{microtype}
  \UseMicrotypeSet[protrusion]{basicmath} % disable protrusion for tt fonts
}{}
\makeatletter
\@ifundefined{KOMAClassName}{% if non-KOMA class
  \IfFileExists{parskip.sty}{%
    \usepackage{parskip}
  }{% else
    \setlength{\parindent}{0pt}
    \setlength{\parskip}{6pt plus 2pt minus 1pt}}
}{% if KOMA class
  \KOMAoptions{parskip=half}}
\makeatother
\usepackage{xcolor}
\usepackage[top=23mm,bottom=23mm,left=24mm,right=24mm]{geometry}
\setlength{\emergencystretch}{3em} % prevent overfull lines
\providecommand{\tightlist}{%
  \setlength{\itemsep}{0pt}\setlength{\parskip}{0pt}}
\setcounter{secnumdepth}{-\maxdimen} % remove section numbering
\usepackage{fontspec}
\usepackage{xeCJK}
\setmainfont{Noto Serif}
\setsansfont{Noto Sans}
\setmonofont{Noto Sans Mono}
\setCJKmainfont[ItalicFont={Noto Serif CJK SC}]{Noto Serif CJK SC}
\setCJKsansfont{Noto Sans CJK SC}
\usepackage{microtype}
\usepackage{setspace}
\setstretch{1.20}
\usepackage[most]{tcolorbox}
\usepackage{fancyhdr}
\usepackage{titlesec}
\usepackage{xcolor}
\usepackage{needspace}
\usepackage{xurl}
\usepackage{enumitem}
\setlist[itemize]{leftmargin=2em,itemsep=3pt,topsep=3pt}
\setlength{\parindent}{1.6em}
\setlength{\parskip}{0.25em}
\raggedbottom
\widowpenalty=10000
\clubpenalty=10000
\displaywidowpenalty=10000
\emergencystretch=3em
\definecolor{primary}{HTML}{2B6259}
\definecolor{muted}{HTML}{546671}
\definecolor{sand}{HTML}{F3F6F5}
\titleformat{\section}{\color{primary}\normalfont\bfseries\Large}{\thesection}{1em}{}
\titleformat{\subsection}{\color{primary}\normalfont\bfseries\large}{\thesubsection}{1em}{}
\titleformat{\subsubsection}{\color{primary}\normalfont\bfseries\normalsize}{\thesubsubsection}{1em}{}
\titlespacing*{\section}{0pt}{2.6ex}{1.2ex}
\titlespacing*{\subsection}{0pt}{2.0ex}{0.8ex}
\titlespacing*{\subsubsection}{0pt}{1.4ex}{0.5ex}
\pagestyle{fancy}
\fancyhf{}
\fancyhead[L]{\small\color{muted} 考研英语（一）\quad 学生讲义}
\fancyhead[R]{\small\color{muted} 语言 · 阅读 · 命题 · 表达}
\fancyfoot[C]{\small\thepage}
\renewcommand{\headrulewidth}{0.35pt}
\renewcommand{\footrulewidth}{0pt}
\setlength{\headheight}{16pt}
\renewenvironment{quote}{\begin{tcolorbox}[enhanced,breakable,colback=sand,colframe=primary!22,boxrule=0.4pt,arc=2pt,left=2.6mm,right=2.6mm,top=1.6mm,bottom=1.6mm]\itshape}{\end{tcolorbox}}
\ifLuaTeX
  \usepackage{selnolig}  % disable illegal ligatures
\fi
\usepackage{bookmark}
\IfFileExists{xurl.sty}{\usepackage{xurl}}{} % add URL line breaks if available
\urlstyle{same}
\hypersetup{
  pdftitle={考研英语（一） · L002},
  pdfauthor={英语能力 · 文章理解 · 考题证据 · 迁移表达},
  hidelinks,
  pdfcreator={LaTeX via pandoc}}

\title{考研英语（一） · L002}
\usepackage{etoolbox}
\makeatletter
\providecommand{\subtitle}[1]{% add subtitle to \maketitle
  \apptocmd{\@title}{\par {\large #1 \par}}{}{}
}
\makeatother
\subtitle{环境与生态｜保护政策与野火管理}
\author{英语能力 · 文章理解 · 考题证据 · 迁移表达}
\date{}

\begin{document}
\maketitle

{
\setcounter{tocdepth}{2}
\tableofcontents
}
\textbf{读懂政策的边界，辨认科学证据的边界}

本册选择两种容易被误读的环境文章。第一篇讨论保护湿地的法律：行政边界与污染流动并不一致。第二篇讨论森林火灾：某些地点在过去几十年里发生的火变少，并不意味着严重野火的风险随之消失。

这两篇并不要求读者成为环境科学家。我们学习的是英语如何说明一个对象的作用、一个规则的适用范围，以及一项证据究竟能支持多强的判断。阅读时先完成自己的判断，再进入讲解和答案部分。

\subsubsection{本册学完以后，你应该能够}\label{ux672cux518cux5b66ux5b8cux4ee5ux540eux4f60ux5e94ux8be5ux80fdux591f}

\begin{itemize}
\item
  能够辨认法律权限、制度范围、跨境影响等抽象词义及搭配。
\item
  能区分火灾发生频率、严重程度和管理难度三种不同命题。
\item
  能在不借助选项的情况下概括作者一段话的作用，并在新文章中复现这种能力。
\item
  能完成有限推断和英译汉，尝试写出有条件、有证据的英文论述。
\end{itemize}

\textbf{使用顺序}：先做导入判断；阅读外部真题并独立答题；使用本册的原文回访任务和语法词汇解释；完成原创阅读与专项训练；最后看答案，再看中英对照附录。完整原卷不在此重印。

\begin{quote}
\textbf{材料身份}　本册出现的外部试题为指定年份的考研英语（一）真题，参考答案在文末注明为``核对后的非官方参考答案''。所有原创例句、迁移阅读和练习均为本册自编，不属于历年真题。
\end{quote}

\section{进入阅读之前｜你能识别这些区别吗}\label{ux8fdbux5165ux9605ux8bfbux4e4bux524dux4f60ux80fdux8bc6ux522bux8fd9ux4e9bux533aux522bux5417}

先读四句英文，不查中文翻译。圈出最可能影响判断的词组，回答后面的中文问题。对照译文统一放在最后。

\textbf{诊断 01}

\begin{quote}
A measure may protect local streams without reducing pollution
downstream.
\end{quote}

这句话否定的是措施的全部价值，还是指出它的作用范围有限？

\begin{center}\rule{0.5\linewidth}{0.5pt}\end{center}

\textbf{诊断 02}

\begin{quote}
The new rule narrowed the agency's authority; it did not abolish the
agency.
\end{quote}

narrow 与 abolish 的差异是什么？

\begin{center}\rule{0.5\linewidth}{0.5pt}\end{center}

\textbf{诊断 03}

\begin{quote}
Although fewer small fires occurred, the forest accumulated more dry
fuel.
\end{quote}

although 将哪两个事实组织成反预期关系？

\begin{center}\rule{0.5\linewidth}{0.5pt}\end{center}

\textbf{诊断 04}

\begin{quote}
Evidence from selected forests cannot establish a universal trend.
\end{quote}

为什么 selected 与 universal 不能随意替换？

\begin{center}\rule{0.5\linewidth}{0.5pt}\end{center}

如果读不顺，先区分障碍：（1）词义不明 （2）句子关系不清
（3）知道句意但不知道为什么作者这样说
（4）文章懂了却不确定选项。下面四个单元会分别处理。

\section{单元一｜湿地：法律可以划界，水却会继续流动}\label{ux5355ux5143ux4e00ux6e7fux5730ux6cd5ux5f8bux53efux4ee5ux5212ux754cux6c34ux5374ux4f1aux7ee7ux7eedux6d41ux52a8}

\textbf{真实试卷}：2024 年英语（一） Reading Text
4（36---40题）。请先打开\href{https://ac-wiki.org/general-skills/study/assets/2024-201.pdf}{可查阅的试卷版本}，标注段号、题号和自己初次选择的答案。

第一遍只看试卷的五个自然段和题目。先写一句自己的概括：这一文章首先写生态作用，为什么随后要讨论法院的判决？不要急于查看答案。

\subsubsection{第一次阅读｜留给自己的证据笔记}\label{ux7b2cux4e00ux6b21ux9605ux8bfbux7559ux7ed9ux81eaux5df1ux7684ux8bc1ux636eux7b14ux8bb0}

请先尝试在纸上完成：

\begin{itemize}
\tightlist
\item
  不看选项，用一句中文写出本文要讨论的问题。
\item
  哪一句/哪一段把``描述对象''转向``提出判断''？
\item
  对最没把握的一题，先在英文原文画出证据，写出自己的英文简答，然后才比较四个选项。
\end{itemize}

\subsubsection{段落阅读实验｜先观察，再返回原文}\label{ux6bb5ux843dux9605ux8bfbux5b9eux9a8cux5148ux89c2ux5bdfux518dux8fd4ux56deux539fux6587}

下列句子是\textbf{原创平行句}，不是从真题抄来的；它们用来演示理解方法。每组请先解释平行句，再回到指定试卷段落完成任务。题目反馈集中在答案部分。

\paragraph{W1 ·
生态作用不是装饰性的开场}\label{w1-ux751fux6001ux4f5cux7528ux4e0dux662fux88c5ux9970ux6027ux7684ux5f00ux573a}

第一段列举湿地对物种、洪水和污染物的作用，读者要分辨的是``功能''而非``景色''。

\textbf{原创平行句}

\begin{quote}
The marsh stores rainwater, filters waste and protects nearby homes.
\end{quote}

\textbf{返回原文的任务}：不看题目选项，用三个英文动词概括 P1
的湿地作用。

\textbf{自己留下的一句话答案}：

\rule{0.75\linewidth}{0.35pt}

\paragraph{W2 ·
裁决改变了什么，没改变什么}\label{w2-ux88c1ux51b3ux6539ux53d8ux4e86ux4ec0ux4e48ux6ca1ux6539ux53d8ux4ec0ux4e48}

第二段从环境转到法律。注意一个名词短语的限定：缩小的是监管权能的适用范围，不是自然资源本身，也不是所有保护措施。

\textbf{原创平行句}

\begin{quote}
The court restricted the board's powers to a narrower group of cases.
\end{quote}

\textbf{返回原文的任务}：找出
P2：法院对联邦监管的适用条件作了什么规定？能否推出``所有湿地均不再受保护''？

\textbf{自己留下的一句话答案}：

\rule{0.75\linewidth}{0.35pt}

\paragraph{W3 ·
例外不能解除跨境风险}\label{w3-ux4f8bux5916ux4e0dux80fdux89e3ux9664ux8de8ux5883ux98ceux9669}

第三段先说本州已有立法，紧接着转折说明仍应担忧。不是作者反对本州法规，而是水体和污染不服从行政辖区。

\textbf{原创平行句}

\begin{quote}
A city may maintain strict standards, yet pollution can arrive from
outside it.
\end{quote}

\textbf{返回原文的任务}：说明P3为何认为``本地有法规''仍不足以让切萨皮克湾高枕无忧。

\textbf{自己留下的一句话答案}：

\rule{0.75\linewidth}{0.35pt}

\paragraph{W4 ·
一个例子为制度协调提供理由}\label{w4-ux4e00ux4e2aux4f8bux5b50ux4e3aux5236ux5ea6ux534fux8c03ux63d0ux4f9bux7406ux7531}

第四段谈联邦机构的协调作用。农民决策与下游捕蟹产业的例子，不是为了指责某个职业，而是说明影响跨行政边界发生。

\textbf{原创平行句}

\begin{quote}
A decision made upstream can impose costs on communities that had no say
in it.
\end{quote}

\textbf{返回原文的任务}：P4中的农业例子支持了哪一项制度性判断？它能否证明所有农民都故意污染？

\textbf{自己留下的一句话答案}：

\rule{0.75\linewidth}{0.35pt}

\paragraph{W5 ·
最后的行动号召回应前文风险}\label{w5-ux6700ux540eux7684ux884cux52a8ux53f7ux53ecux56deux5e94ux524dux6587ux98ceux9669}

第五段是社论的倡议，而不是新的实验结果。辨认作者在呼吁谁、做什么、为什么。

\textbf{原创平行句}

\begin{quote}
The editorial urges lawmakers to review safeguards before the damage
spreads.
\end{quote}

\textbf{返回原文的任务}：P5的对象是州立法者还是环保科学家？作者呼吁审视保护措施还是仅扩大避难所？

\textbf{自己留下的一句话答案}：

\rule{0.75\linewidth}{0.35pt}

\subsubsection{词汇与表达｜从具体误解入手}\label{ux8bcdux6c47ux4e0eux8868ux8fbeux4eceux5177ux4f53ux8befux89e3ux5165ux624b}

这不是统一格式的``词义抄写表''。请先理解词在本文中的意义，关注它改变了哪一层断言，再试着不看中文译出原创例句。无需把所有义项一口气背完。

\paragraph{01　filter}\label{filter}

\textbf{本篇核心用法}　v. 过滤；除去。　\textbf{典型组合}　filter
pollutants / filter out impurities。

\textbf{容易误解的地方}　强调通过介质去除某物；与``筛选信息''的抽象义相通。

\textbf{原创语境句}

\begin{quote}
Wetlands can filter some pollutants before water enters the reservoir.
\end{quote}

\textbf{再想一步}　这个句子能否证明湿地可以消除一切污染？

\paragraph{02　authority}\label{authority}

\textbf{本篇核心用法}　n.~法定权力；权限。　\textbf{典型组合}　regulatory
authority / have authority to do。

\textbf{容易误解的地方}　在P2是``有权做某事''，不是某位``权威专家''。

\textbf{原创语境句}

\begin{quote}
Local councils lack the authority to change national safety standards.
\end{quote}

\textbf{再想一步}　lack authority 与 choose not to act
哪个表达权限缺失？

\paragraph{03　scope}\label{scope}

\textbf{本篇核心用法}　n.~范围；适用边界。　\textbf{典型组合}　the scope
of a rule / narrow the scope。

\textbf{容易误解的地方}　规则可以变窄而继续存在；不要译成``规模''后忽略其适用范围。

\textbf{原创语境句}

\begin{quote}
The amendment reduced the scope of the inspection programme.
\end{quote}

\textbf{再想一步}　减少检查计划覆盖面等于彻底停止检查吗？

\paragraph{04　at odds with}\label{at-odds-with}

\textbf{本篇核心用法}　与\ldots\ldots 不一致；发生冲突。　\textbf{典型组合}　be
at odds with regulations / evidence。

\textbf{容易误解的地方}　表示利益或判断不相容，不自动断定违法。

\textbf{原创语境句}

\begin{quote}
The company's expansion plans were at odds with the town's conservation
goals.
\end{quote}

\textbf{再想一步}　这句话断言扩张计划一定违法了吗？

\paragraph{05　repercussion}\label{repercussion}

\textbf{本篇核心用法}　n.~连带后果（通常负面）。　\textbf{典型组合}　serious
repercussions / long-term repercussions。

\textbf{容易误解的地方}　强调后续波及的影响，常用复数；可与 immediate
effect 对照。

\textbf{原创语境句}

\begin{quote}
The abrupt change had repercussions for farms far beyond the city.
\end{quote}

\textbf{再想一步}　这里为什么用 repercussion 而不只说 change？

\paragraph{06　jurisdiction}\label{jurisdiction}

\textbf{本篇核心用法}　n.~管辖权；司法或行政管辖区域。　\textbf{典型组合}　within
a jurisdiction / cross jurisdictions。

\textbf{容易误解的地方}　与 authority
不同：前者强调``在哪里、对什么事项有管辖''，后者强调``能依法做什么''。

\textbf{原创语境句}

\begin{quote}
The river passes through several jurisdictions with different rules.
\end{quote}

\textbf{再想一步}　一个流域跨越多个辖区，会给统一治理带来什么问题？

\paragraph{07　shortsighted}\label{shortsighted}

\textbf{本篇核心用法}　adj. 缺乏长远眼光的。　\textbf{典型组合}　a
shortsighted policy / view。

\textbf{容易误解的地方}　是作者的负面评价，并非``近视''这一生理义。

\textbf{原创语境句}

\begin{quote}
Cutting flood prevention funds may be shortsighted even if it saves
money today.
\end{quote}

\textbf{再想一步}　如何区分作者的评价与一项客观的财政事实？

\paragraph{08　transcend}\label{transcend}

\textbf{本篇核心用法}　v. 超越；跨越。　\textbf{典型组合}　transcend
borders / divisions。

\textbf{容易误解的地方}　对象可以是分歧、局部利益或限制，并不意味着消除所有差异。

\textbf{原创语境句}

\begin{quote}
Cooperation can transcend state borders when rivers connect communities.
\end{quote}

\textbf{再想一步}　合作跨界等于州界在法律上不存在了吗？

\paragraph{09　runoff}\label{runoff}

\textbf{本篇核心用法}　n.~地表径流。　\textbf{典型组合}　agricultural
runoff / stormwater runoff。

\textbf{容易误解的地方}　指流过土地表面并进入水体的水，可能夹带污染物。

\textbf{原创语境句}

\begin{quote}
Heavy rainfall carried fertilizer runoff into the nearby lake.
\end{quote}

\textbf{再想一步}　句中污染物通过哪一媒介移动？

\paragraph{10　call on}\label{call-on}

\textbf{本篇核心用法}　v. phr. 呼吁；要求。　\textbf{典型组合}　call on
lawmakers to review。

\textbf{容易误解的地方}　此处是公开要求采取行动，不是``拜访某人''。

\textbf{原创语境句}

\begin{quote}
Residents called on the council to publish its water-quality data.
\end{quote}

\textbf{再想一步}　call on 后的行动由谁实施？

\paragraph{11　stakes}\label{stakes}

\textbf{本篇核心用法}　n.~利害关系；风险所系。　\textbf{典型组合}　what
is at stake / high stakes。

\textbf{容易误解的地方}　所涉及的利益可能是生态、生活或长期经济损失。

\textbf{原创语境句}

\begin{quote}
The debate involves high stakes for fishing towns downstream.
\end{quote}

\textbf{再想一步}　stake 指的是风险所系，而不是赌局吗？

\subsubsection{与上一册连接的旧词}\label{ux4e0eux4e0aux4e00ux518cux8fdeux63a5ux7684ux65e7ux8bcd}

\begin{itemize}
\tightlist
\item
  regulate（S001）---监管与规则适用，并非简单``控制一切''。
\end{itemize}

\subsubsection{回到试卷：给初次判断一次修正机会}\label{ux56deux5230ux8bd5ux5377ux7ed9ux521dux6b21ux5224ux65adux4e00ux6b21ux4feeux6b63ux673aux4f1a}

再做一遍这一篇真题的五道题；不要直接照抄第一次答案。对每道题写出：\textbf{原文哪一处提供证据？题干要求概括事实、判断态度还是进行推断？最诱人的错误项改变了什么？}

\begin{itemize}
\item
  \textbf{第36题}　证据段：\_\_\_\_　简答另写：\_\_\_\_　最易误选项：\_\_\_\_
\item
  \textbf{第37题}　证据段：\_\_\_\_　简答另写：\_\_\_\_　最易误选项：\_\_\_\_
\item
  \textbf{第38题}　证据段：\_\_\_\_　简答另写：\_\_\_\_　最易误选项：\_\_\_\_
\item
  \textbf{第39题}　证据段：\_\_\_\_　简答另写：\_\_\_\_　最易误选项：\_\_\_\_
\item
  \textbf{第40题}　证据段：\_\_\_\_　简答另写：\_\_\_\_　最易误选项：\_\_\_\_
\end{itemize}

\section{单元二｜野火：少发生的火，为什么可能留下更多风险}\label{ux5355ux5143ux4e8cux91ceux706bux5c11ux53d1ux751fux7684ux706bux4e3aux4ec0ux4e48ux53efux80fdux7559ux4e0bux66f4ux591aux98ceux9669}

\textbf{真实试卷}：2026 年英语（一） Reading Text
4（36---40题）。请先打开\href{https://english-exam.lazynote.cn/kaoyan/sections/2026-english-one/section2-part-a-4/}{可查阅的试卷版本}，标注段号、题号和自己初次选择的答案。

先分别写下三个英文名词：frequency（频率）、severity（严重程度）、risk（风险）。读完以后检查，文章在哪几段谈到的实际上是不同量？

\subsubsection{第一次阅读｜留给自己的证据笔记}\label{ux7b2cux4e00ux6b21ux9605ux8bfbux7559ux7ed9ux81eaux5df1ux7684ux8bc1ux636eux7b14ux8bb0-1}

请先尝试在纸上完成：

\begin{itemize}
\tightlist
\item
  不看选项，用一句中文写出本文要讨论的问题。
\item
  哪一句/哪一段把``描述对象''转向``提出判断''？
\item
  对最没把握的一题，先在英文原文画出证据，写出自己的英文简答，然后才比较四个选项。
\end{itemize}

\subsubsection{段落阅读实验｜先观察，再返回原文}\label{ux6bb5ux843dux9605ux8bfbux5b9eux9a8cux5148ux89c2ux5bdfux518dux8fd4ux56deux539fux6587-1}

下列句子是\textbf{原创平行句}，不是从真题抄来的；它们用来演示理解方法。每组请先解释平行句，再回到指定试卷段落完成任务。题目反馈集中在答案部分。

\paragraph{F1 ·
个人经验遇上历史证据}\label{f1-ux4e2aux4ebaux7ecfux9a8cux9047ux4e0aux5386ux53f2ux8bc1ux636e}

第一段从Swetnam早年的灭火职责进入树轮证据。惊讶的对象是殖民前火灾留下的多次痕迹，不是70年代的消防技术。

\textbf{原创平行句}

\begin{quote}
Old trees preserve evidence that a recent survey may overlook.
\end{quote}

\textbf{返回原文的任务}：指出第36题要求找的是人物``发现了什么''，还是``任职时的工作目标''。

\textbf{自己留下的一句话答案}：

\rule{0.75\linewidth}{0.35pt}

\paragraph{F2 ·
低烈度火灾的功能}\label{f2-ux4f4eux70c8ux5ea6ux706bux707eux7684ux529fux80fd}

第二段不是讨论如何扑灭每一场火，而是说明低强度燃烧怎样清理落叶、促进生长和减少病害。

\textbf{原创平行句}

\begin{quote}
A mild burn can remove dry litter without destroying mature trees.
\end{quote}

\textbf{返回原文的任务}：为什么 P2
的主旨应概括为``多种好处''，而不只聚焦某一种病树？

\textbf{自己留下的一句话答案}：

\rule{0.75\linewidth}{0.35pt}

\paragraph{F3 ·
观察指标和比较基线}\label{f3-ux89c2ux5bdfux6307ux6807ux548cux6bd4ux8f83ux57faux7ebf}

第三段的数据讨论具体研究点位和历史记录，对比的是特定时期的火灾发生次数。不能偷换为全球火灾危害正在全面下降。

\textbf{原创平行句}

\begin{quote}
At the sampled sites, recent fire counts were lower than the historical
pattern suggested.
\end{quote}

\textbf{返回原文的任务}：研究结论的地域/指标/比较对象各是什么？

\textbf{自己留下的一句话答案}：

\rule{0.75\linewidth}{0.35pt}

\paragraph{F4 ·
表面成功与累积性代价}\label{f4-ux8868ux9762ux6210ux529fux4e0eux7d2fux79efux6027ux4ee3ux4ef7}

第四段既承认现代灭火的成就，又指出持续压制燃烧导致林下可燃物累积，与气候变化共同使极端火情更难处理。

\textbf{原创平行句}

\begin{quote}
Suppressing every small fire may allow dry fuel to accumulate over time.
\end{quote}

\textbf{返回原文的任务}：为什么``曾经有效''与``长期可能增加严重风险''并不矛盾？

\textbf{自己留下的一句话答案}：

\rule{0.75\linewidth}{0.35pt}

\paragraph{F5 ·
有计划的燃烧与条件限定}\label{f5-ux6709ux8ba1ux5212ux7684ux71c3ux70e7ux4e0eux6761ux4ef6ux9650ux5b9a}

第五段提出按天气及地形条件精心安排的火；prescribed不是``医生开药''，是有组织的生态管理措施。

\textbf{原创平行句}

\begin{quote}
A prescribed burn is planned for conditions in which the flames can be
managed.
\end{quote}

\textbf{返回原文的任务}：作者称其有用，是否就主张在任何天气下点火？

\textbf{自己留下的一句话答案}：

\rule{0.75\linewidth}{0.35pt}

\paragraph{F6 ·
双重困难与有限共识}\label{f6-ux53ccux91cdux56f0ux96beux4e0eux6709ux9650ux5171ux8bc6}

末段的双刃剑：气候变暖增加对计划用火的需求，同时缩小了安全实施窗口。两位专家共识是不能忽视有计划用火。

\textbf{原创平行句}

\begin{quote}
The policy may become more necessary just as it becomes harder to carry
out safely.
\end{quote}

\textbf{返回原文的任务}：两位专家共同支持什么？谁指出了气候变化导致的实施困难？

\textbf{自己留下的一句话答案}：

\rule{0.75\linewidth}{0.35pt}

\subsubsection{词汇与表达｜从具体误解入手}\label{ux8bcdux6c47ux4e0eux8868ux8fbeux4eceux5177ux4f53ux8befux89e3ux5165ux624b-1}

这不是统一格式的``词义抄写表''。请先理解词在本文中的意义，关注它改变了哪一层断言，再试着不看中文译出原创例句。无需把所有义项一口气背完。

\paragraph{01　mandate}\label{mandate}

\textbf{本篇核心用法}　n.~任务；授权职责。　\textbf{典型组合}　a mandate
to protect / under a mandate。

\textbf{容易误解的地方}　这里是机构或个人承担的职责，不是药品``命令''。

\textbf{原创语境句}

\begin{quote}
Her mandate was to reduce immediate threats to nearby villages.
\end{quote}

\textbf{再想一步}　职责是否已经说明履职方法永远正确？

\paragraph{02　curb}\label{curb}

\textbf{本篇核心用法}　v. 抑制；约束。　\textbf{典型组合}　curb the
spread / curb emissions。

\textbf{容易误解的地方}　通常是限制规模与速度，不保证清零。

\textbf{原创语境句}

\begin{quote}
New measures have curbed the spread of small grass fires.
\end{quote}

\textbf{再想一步}　curb 与 eliminate 的断言强度有什么差别？

\paragraph{03　peter out}\label{peter-out}

\textbf{本篇核心用法}　phr. 逐渐减弱直至停止。　\textbf{典型组合}　the
fire petered out。

\textbf{容易误解的地方}　强调缓慢消退的过程，常用于火、声音、计划等。

\textbf{原创语境句}

\begin{quote}
The small fire petered out as the weather became cooler.
\end{quote}

\textbf{再想一步}　这是否等同于消防人员立即扑灭？

\paragraph{04　compile}\label{compile}

\textbf{本篇核心用法}　v. 汇编；整理。　\textbf{典型组合}　compile
records / compile a dataset。

\textbf{容易误解的地方}　强调从已有资料收集整理，不意味着``编造''。

\textbf{原创语境句}

\begin{quote}
Scientists compiled records from trees across the region.
\end{quote}

\textbf{再想一步}　汇编数据不等于创造新事实。

\paragraph{05　shortfall}\label{shortfall}

\textbf{本篇核心用法}　n.~不足；缺口。　\textbf{典型组合}　a funding
shortfall / fire shortfall。

\textbf{容易误解的地方}　须先说明与什么预期或需求相比``缺少''。

\textbf{原创语境句}

\begin{quote}
The survey found a shortfall in routine maintenance compared with the
previous decade.
\end{quote}

\textbf{再想一步}　没有基准线，能否单凭 shortfall 判断变化程度？

\paragraph{06　suppression}\label{suppression}

\textbf{本篇核心用法}　n.~压制；扑灭。　\textbf{典型组合}　fire
suppression / suppression policy。

\textbf{容易误解的地方}　fire suppression
指控制或扑灭野火，是行动或政策，不是发生频率。

\textbf{原创语境句}

\begin{quote}
Decades of suppression altered the conditions under the forest canopy.
\end{quote}

\textbf{再想一步}　为什么行动的成功不等于长期生态结果只有好处？

\paragraph{07　testament}\label{testament}

\textbf{本篇核心用法}　n.~有力证明；体现。　\textbf{典型组合}　a
testament to skill / effort。

\textbf{容易误解的地方}　常在 a testament to 结构中给出某成就的意义。

\textbf{原创语境句}

\begin{quote}
The restored trail is a testament to the volunteers' patience.
\end{quote}

\textbf{再想一步}　这比 report a result 多出了何种评价色彩？

\paragraph{08　erode}\label{erode}

\textbf{本篇核心用法}　v. 侵蚀；逐渐削弱。　\textbf{典型组合}　erode
trust / erode capacity。

\textbf{容易误解的地方}　抽象用法强调逐步削弱能力，不等于完全失去。

\textbf{原创语境句}

\begin{quote}
Repeated delays eroded residents' trust in the warning system.
\end{quote}

\textbf{再想一步}　eroded trust 能否推出居民已经毫无信任？

\paragraph{09　ignite}\label{ignite}

\textbf{本篇核心用法}　v. 点燃；引发。　\textbf{典型组合}　ignite a fire
/ ignite debate。

\textbf{容易误解的地方}　既可指真实点燃，也能用于抽象争论。

\textbf{原创语境句}

\begin{quote}
A single spark can ignite a fire in unusually dry grass.
\end{quote}

\textbf{再想一步}　can 表示可能性还是必然性？

\paragraph{10　prescribed}\label{prescribed}

\textbf{本篇核心用法}　adj.
按规定计划的；预先安排的。　\textbf{典型组合}　prescribed burns /
prescribed conditions。

\textbf{容易误解的地方}　在火管理中指预定条件下人为施加燃烧；不等同于药物处方。

\textbf{原创语境句}

\begin{quote}
The team postponed the prescribed burn because the wind shifted.
\end{quote}

\textbf{再想一步}　推迟是否说明该措施从原理上无用？

\paragraph{11　tactic}\label{tactic}

\textbf{本篇核心用法}　n.~具体手段；策略。　\textbf{典型组合}　a
management tactic / use a tactic。

\textbf{容易误解的地方}　比 strategy
更偏向具体行动，但两词实际使用常有交叉。

\textbf{原创语境句}

\begin{quote}
Managers changed their tactic when the weather became unpredictable.
\end{quote}

\textbf{再想一步}　策略因条件变化而调整，是否意味着目标改变？

\paragraph{12　double-edged sword}\label{double-edged-sword}

\textbf{本篇核心用法}　n.~两面都有后果的事物。　\textbf{典型组合}　a
double-edged sword。

\textbf{容易误解的地方}　表达同一变化带来彼此牵制的收益与成本，不是``结果无法判断''。

\textbf{原创语境句}

\begin{quote}
Improved access is a double-edged sword: it supports tourism but
increases pressure on habitats.
\end{quote}

\textbf{再想一步}　请说明两个方向分别影响什么。

\subsubsection{与上一册连接的旧词}\label{ux4e0eux4e0aux4e00ux518cux8fdeux63a5ux7684ux65e7ux8bcd-1}

\begin{itemize}
\item
  resilient（L001）---能经受冲击并恢复，不等于永远不受损。
\item
  subsequent（L001）---发生在之后，不能直接推出由前事件引起。
\end{itemize}

\subsubsection{回到试卷：给初次判断一次修正机会}\label{ux56deux5230ux8bd5ux5377ux7ed9ux521dux6b21ux5224ux65adux4e00ux6b21ux4feeux6b63ux673aux4f1a-1}

再做一遍这一篇真题的五道题；不要直接照抄第一次答案。对每道题写出：\textbf{原文哪一处提供证据？题干要求概括事实、判断态度还是进行推断？最诱人的错误项改变了什么？}

\begin{itemize}
\item
  \textbf{第36题}　证据段：\_\_\_\_　简答另写：\_\_\_\_　最易误选项：\_\_\_\_
\item
  \textbf{第37题}　证据段：\_\_\_\_　简答另写：\_\_\_\_　最易误选项：\_\_\_\_
\item
  \textbf{第38题}　证据段：\_\_\_\_　简答另写：\_\_\_\_　最易误选项：\_\_\_\_
\item
  \textbf{第39题}　证据段：\_\_\_\_　简答另写：\_\_\_\_　最易误选项：\_\_\_\_
\item
  \textbf{第40题}　证据段：\_\_\_\_　简答另写：\_\_\_\_　最易误选项：\_\_\_\_
\end{itemize}

\section{句子工坊｜精准理解比翻译顺口更重要}\label{ux53e5ux5b50ux5de5ux574aux7cbeux51c6ux7406ux89e3ux6bd4ux7ffbux8bd1ux987aux53e3ux66f4ux91cdux8981}

下面六句均为原创，请先划分主干，再圈出限定词、连接词或修饰范围。\textbf{先用英语或中文说清命题，再翻译}；不要直接猜中文意思。

\subsubsection{S1｜句子观察}\label{s1ux53e5ux5b50ux89c2ux5bdf}

\begin{quote}
When the boundary of a rule changes, the river itself remains unchanged.
\end{quote}

\textbf{观察任务}　指出when引导的时间/条件背景；找出主句反驳的误解。

\textbf{独立翻译}　

\rule{0.8\linewidth}{0.35pt}

\subsubsection{S2｜句子观察}\label{s2ux53e5ux5b50ux89c2ux5bdf}

\begin{quote}
The state had adopted strict controls, yet it could not prevent
pollution arriving from upstream.
\end{quote}

\textbf{观察任务}　解释yet连接的两个命题以及它为何重要。

\textbf{独立翻译}　

\rule{0.8\linewidth}{0.35pt}

\subsubsection{S3｜句子观察}\label{s3ux53e5ux5b50ux89c2ux5bdf}

\begin{quote}
The evidence suggests a decline in fire frequency at these sites, not a
universal decline in fire danger.
\end{quote}

\textbf{观察任务}　划出结论的范围、指标及被否定的过度推断。

\textbf{独立翻译}　

\rule{0.8\linewidth}{0.35pt}

\subsubsection{S4｜句子观察}\label{s4ux53e5ux5b50ux89c2ux5bdf}

\begin{quote}
It is not the number of interventions alone but their long-term effects
that should guide policy.
\end{quote}

\textbf{观察任务}　识别not\ldots alone but\ldots 的对比焦点。

\textbf{独立翻译}　

\rule{0.8\linewidth}{0.35pt}

\subsubsection{S5｜句子观察}\label{s5ux53e5ux5b50ux89c2ux5bdf}

\begin{quote}
Because the planned fire was delayed, the team had time to check whether
conditions were safe.
\end{quote}

\textbf{观察任务}　区分because因果与whether从句表达的不确定判断。

\textbf{独立翻译}　

\rule{0.8\linewidth}{0.35pt}

\subsubsection{S6｜句子观察}\label{s6ux53e5ux5b50ux89c2ux5bdf}

\begin{quote}
Even if an intervention reduces immediate losses, it may create risks
that appear only years later.
\end{quote}

\textbf{观察任务}　找出让步前提、may的力度和that定语从句的先行词。

\textbf{独立翻译}　

\rule{0.8\linewidth}{0.35pt}

\textbf{读者提示}：如果你能流畅译出一个句子，却无法说明``谁在什么条件下断言了什么''，仍需要回到连词、限定和先行词。真正重要的是恢复命题关系。

\section{阅读迁移｜两篇没有见过的英文文章}\label{ux9605ux8bfbux8fc1ux79fbux4e24ux7bc7ux6ca1ux6709ux89c1ux8fc7ux7684ux82f1ux6587ux6587ux7ae0}

以下均为\textbf{原创训练文章}。先按真实阅读的方式独立完成五道选择题和两道无选项简答，再读答案。这里的目标不是猜同一个题型，而是检验词义、篇章和证据判断能否迁移到新题材。

\subsection{Reading A · The Quiet Value of a
Floodplain}\label{reading-a-the-quiet-value-of-a-floodplain}

原创迁移阅读，主题为城镇防洪与洪泛区，不是2024年真题改写。　\textbf{建议首读限时：7---9分钟；不要先翻到附录看中文。}

\textbf{P1}　A town on the River Lorne spent years strengthening the
wall that separated its shopping district from the water. The wall
worked during ordinary storms. Shops stayed open, and local officials
described the project as a success. But after an unusually wet winter,
engineers discovered that the wall had pushed more water towards
villages downstream. Protection in one place had changed the
distribution of danger, rather than removed it.

\textbf{P2}　The regional council then considered restoring a floodplain
north of the town. Under the proposal, part of the farmland would be
allowed to flood during periods of heavy rain. Supporters argued that
temporary flooding on open land could prevent far greater damage in
densely populated areas. Farmers, however, pointed out that even
temporary losses could place serious pressure on families whose income
depended on a single season.

\textbf{P3}　Neither side disputed the need for flood protection. Their
disagreement concerned who should bear its costs and how those costs
should be measured. If the council counted only repairs to buildings,
the floodplain looked inexpensive. If it included lost crops, disrupted
transport and the uncertainty faced by farmers, the calculation became
more complicated.

\textbf{P4}　The final plan combined smaller walls with a compensation
fund for affected farms. It also required the council to publish annual
reports on flooding across the whole river basin. The arrangement did
not guarantee that every household would be safe. It did, however, make
the trade-offs visible and allow communities to judge whether protection
was being shared fairly.

\subsubsection{Questions ·
独立答题}\label{questions-ux72ecux7acbux7b54ux9898}

\textbf{A1. What did the engineers discover after the wet winter?}

\begin{itemize}
\item
  \textbf{A.} The wall had failed in every storm.
\item
  \textbf{B.} The wall had shifted some risk downstream.
\item
  \textbf{C.} The town had no need for flood protection.
\item
  \textbf{D.} The villages had built a second wall.
\end{itemize}

\textbf{A2. The farmers were mainly concerned about}

\begin{itemize}
\item
  \textbf{A.} the cost of replacing the town's shops.
\item
  \textbf{B.} the absence of an annual report.
\item
  \textbf{C.} the pressure of seasonal losses on family income.
\item
  \textbf{D.} the closing of the regional council.
\end{itemize}

\textbf{A3. The third paragraph mainly examines}

\begin{itemize}
\item
  \textbf{A.} how costs are defined and counted.
\item
  \textbf{B.} why farmers oppose all public spending.
\item
  \textbf{C.} whether shop repairs are legally required.
\item
  \textbf{D.} when the river changed its course.
\end{itemize}

\textbf{A4. The compensation fund was intended to}

\begin{itemize}
\item
  \textbf{A.} eliminate the possibility of flooding.
\item
  \textbf{B.} replace all flood reports.
\item
  \textbf{C.} guarantee profits to every business.
\item
  \textbf{D.} address losses borne by affected farms.
\end{itemize}

\textbf{A5. Which statement best captures the author's point?}

\begin{itemize}
\item
  \textbf{A.} Good protection requires attention to where costs fall.
\item
  \textbf{B.} Local walls should never be built.
\item
  \textbf{C.} Farmers should determine all public policy.
\item
  \textbf{D.} Annual reports can prevent natural disasters.
\end{itemize}

\subsubsection{No-option questions ·
不依赖选项的理解}\label{no-option-questions-ux4e0dux4f9dux8d56ux9009ux9879ux7684ux7406ux89e3}

\textbf{AO1.} Why does the author call the final plan a trade-off?

回答时至少指出原文依据及必要限定。

\textbf{AO2.} Does publication of reports itself ensure fairness?

回答时至少指出原文依据及必要限定。

\subsubsection{读后复盘}\label{ux8bfbux540eux590dux76d8}

闭卷用约40字写出：作者起初介绍的问题、最重要的转向、以及末段的有限结论。然后回到文中画出支撑这三点的句子。

\subsection{Reading B · The Fires a Park Did Not
Have}\label{reading-b-the-fires-a-park-did-not-have}

原创迁移阅读，讨论公园小火管理与监测指标，非2026年原文改写。　\textbf{建议首读限时：7---9分钟；不要先翻到附录看中文。}

\textbf{P1}　For decades, staff at Alder Park judged their fire
programme by a simple number: how few fires occurred each summer. The
fewer incidents they recorded, the more successful the programme
appeared. Visitors felt reassured, and managers could report that fewer
people had been exposed to smoke. Yet the measure offered no information
about what was happening to the forest floor.

\textbf{P2}　A survey later found thick layers of dry branches in areas
where small burns had long been suppressed. Ecologists did not argue
that every wildfire was beneficial. They distinguished mild fires, which
might clear some fallen material, from extreme fires capable of damaging
mature trees and nearby buildings. Treating both as a single category
had hidden a difference important to future decisions.

\textbf{P3}　The park proposed a limited trial of planned burning. Fires
would be introduced only on selected days, after checks of wind,
moisture and access roads. A small team would monitor the operation, and
nearby residents would be informed in advance. Critics reasonably asked
whether changing weather could turn a controlled burn into an emergency.
The project team acknowledged that this risk could not be reduced to
zero.

\textbf{P4}　For this reason, the trial would be evaluated by several
measures rather than a single fire count: fuel levels, damage from later
fires, costs of preparation and any harm caused by the trial itself. The
aim was not to celebrate fire. It was to learn which combination of
prevention and carefully managed burning could make the park safer over
time.

\subsubsection{Questions ·
独立答题}\label{questions-ux72ecux7acbux7b54ux9898-1}

\textbf{B1. The park originally assessed success mainly through}

\begin{itemize}
\item
  \textbf{A.} the number of summer fires.
\item
  \textbf{B.} the age of its oldest trees.
\item
  \textbf{C.} the amount of rain in winter.
\item
  \textbf{D.} the cost of public transport.
\end{itemize}

\textbf{B2. The ecologists' distinction in Paragraph 2 suggests that}

\begin{itemize}
\item
  \textbf{A.} all fires should burn freely.
\item
  \textbf{B.} fire severity affects how an event is assessed.
\item
  \textbf{C.} smoke is always harmless.
\item
  \textbf{D.} old trees prevent extreme fires.
\end{itemize}

\textbf{B3. Why would the trial be limited to selected days?}

\begin{itemize}
\item
  \textbf{A.} To maximise the total number of fires.
\item
  \textbf{B.} To avoid telling residents.
\item
  \textbf{C.} To test risks under assessed conditions.
\item
  \textbf{D.} To stop all forest surveys.
\end{itemize}

\textbf{B4. The project team conceded that planned burning}

\begin{itemize}
\item
  \textbf{A.} had already destroyed every building.
\item
  \textbf{B.} would never be useful.
\item
  \textbf{C.} required no monitoring.
\item
  \textbf{D.} could not be made completely risk-free.
\end{itemize}

\textbf{B5. The central argument is that managers should}

\begin{itemize}
\item
  \textbf{A.} replace all prevention with fires.
\item
  \textbf{B.} judge measures by several long-term outcomes.
\item
  \textbf{C.} avoid listening to local residents.
\item
  \textbf{D.} prefer immediate numbers to consequences.
\end{itemize}

\subsubsection{No-option questions ·
不依赖选项的理解}\label{no-option-questions-ux4e0dux4f9dux8d56ux9009ux9879ux7684ux7406ux89e3-1}

\textbf{BO1.} What is the weakness of counting fires alone?

回答时至少指出原文依据及必要限定。

\textbf{BO2.} Does the author recommend uncontrolled wildfires?

回答时至少指出原文依据及必要限定。

\subsubsection{读后复盘}\label{ux8bfbux540eux590dux76d8-1}

闭卷用约40字写出：作者起初介绍的问题、最重要的转向、以及末段的有限结论。然后回到文中画出支撑这三点的句子。

\section{考试专项衔接｜不只训练四选一阅读}\label{ux8003ux8bd5ux4e13ux9879ux8854ux63a5ux4e0dux53eaux8badux7ec3ux56dbux9009ux4e00ux9605ux8bfb}

以下练习把本册词句带入词汇使用、语篇衔接、翻译和写作。题目均为原创，先完成再看答案。

\subsection{A　语境选词｜结构与词义同时成立}\label{a-ux8bedux5883ux9009ux8bcdux7ed3ux6784ux4e0eux8bcdux4e49ux540cux65f6ux6210ux7acb}

\textbf{词库（无需全部用完）}：narrow, runoff, jurisdictions, compiled,
erode, call on, prescribed, double-edged
sword。必要时改变动词形态或填入短语的适当形式。

\textbf{V01.} The proposal would \_\_\_\_\_ the agency's powers to a
smaller set of cases.

\textbf{V02.} Rainwater carried agricultural \_\_\_\_\_ into the river
after the storm.

\textbf{V03.} The river crosses several \_\_\_\_\_, each with its own
rules.

\textbf{V04.} Scientists \_\_\_\_\_ fire records collected from hundreds
of trees.

\textbf{V05.} A series of small cuts may gradually \_\_\_\_\_ public
trust.

\textbf{V06.} The council decided to \_\_\_\_\_ residents to take part
in the review.

\textbf{V07.} The carefully \_\_\_\_\_ fire was delayed because the wind
changed.

\textbf{V08.} A rule can be a \_\_\_\_\_ when it protects one group but
burdens another.

\subsection{B　篇章衔接｜从段落功能而不是重复词作答}\label{b-ux7bc7ux7ae0ux8854ux63a5ux4eceux6bb5ux843dux529fux80fdux800cux4e0dux662fux91cdux590dux8bcdux4f5cux7b54}

\textbf{篇章衔接：同一条河，两种政策目标}。根据 G1、G2 前后的信息，从
A---D 选择适当的句子；每项最多使用一次。

The committee first measured how much rainwater each district could
store. It reported that the largest upstream wetland could hold a
significant amount during heavy storms.

{[}G1{]}

Officials then discovered that local budgets were prepared separately,
even though flooding crossed district boundaries.

{[}G2{]}

The new agreement therefore shared monitoring costs and required annual
public reports.

\begin{itemize}
\item
  \textbf{A.} Yet this storage capacity did not explain who should pay
  for its protection.
\item
  \textbf{B.} The scientists refused to measure any rainfall at all.
\item
  \textbf{C.} Without coordination, no district could be sure how its
  neighbours would act.
\item
  \textbf{D.} The upstream wetland was built by the committee only
  yesterday.
\end{itemize}

\subsection{C　英译汉｜信息不能丢，也不能增加}\label{c-ux82f1ux8bd1ux6c49ux4fe1ux606fux4e0dux80fdux4e22ux4e5fux4e0dux80fdux589eux52a0}

\textbf{T01.} Although the county improved its own flood barriers,
pollution continued to arrive from areas outside its jurisdiction.

译文：

\rule{0.8\linewidth}{0.35pt}

\textbf{T02.} The evidence concerns the frequency of small fires at
selected sites rather than the severity of every wildfire.

译文：

\rule{0.8\linewidth}{0.35pt}

\textbf{T03.} A method that reduces immediate danger may nevertheless
create a long-term risk if it changes the conditions of the forest.

译文：

\rule{0.8\linewidth}{0.35pt}

\textbf{T04.} Instead of asking whether a policy is good or bad in
general, we should examine which groups benefit and which bear the
costs.

译文：

\rule{0.8\linewidth}{0.35pt}

\subsection{D　汉译英｜使用准确而自然的表达}\label{d-ux6c49ux8bd1ux82f1ux4f7fux7528ux51c6ux786eux800cux81eaux7136ux7684ux8868ux8fbe}

\textbf{P01.} 必须说明干预措施的适用范围，不能把局部成功推广到全部流域。

试译：

\rule{0.8\linewidth}{0.35pt}

\textbf{P02.}
计划用火有助于减少某些风险，但其实施必须依赖适当的天气条件。

试译：

\rule{0.8\linewidth}{0.35pt}

\subsection{E　写作小任务｜把语言知识用于表达}\label{e-ux5199ux4f5cux5c0fux4efbux52a1ux628aux8bedux8a00ux77e5ux8bc6ux7528ux4e8eux8868ux8fbe}

\textbf{任务（英文）}　Write about 90--120 words for a school
newsletter. Explain why one environmental indicator may be insufficient
for judging the success of a policy. Use an example and suggest a fairer
way of assessment.

\textbf{完成要求}　先拟两到三个论点，每段一句主题句。必须有一个明确的限定（some
/ may / under certain
conditions等），并说明例子支持的范围。写完自查：是否把一种可能夸大为必然？

\textbf{自己的草稿}

\vspace{2cm}\rule{0.95\linewidth}{0.35pt}\vspace{0.6cm}

\clearpage

\section{答案、证据与思维诊断}\label{ux7b54ux6848ux8bc1ux636eux4e0eux601dux7ef4ux8bcaux65ad}

这里既不是``答案字母库''，也不是``命题者的秘密''。重点是逐步验证自己的理解：\textbf{词义成立吗、原文说了吗、题干要求这一层信息吗、错误项改变了什么？}

\subsection{一　开始前的四道判断}\label{ux4e00-ux5f00ux59cbux524dux7684ux56dbux9053ux5224ux65ad}

\textbf{诊断 01}　区分局部效益与全流域结果，不能以 without
后的限定否定整项措施。

\textbf{诊断 02}　narrow 是收窄权限，abolish 是撤销实体，两者不同。

\textbf{诊断 03}　让步连接了数量与积累风险方向不同的事实。

\textbf{诊断 04}　selected 把证据限定在样本范围，不能推广所有地点。

\subsection{二　原卷段落回访｜按证据回答，不按印象记忆}\label{ux4e8c-ux539fux5377ux6bb5ux843dux56deux8bbfux6309ux8bc1ux636eux56deux7b54ux4e0dux6309ux5370ux8c61ux8bb0ux5fc6}

\subsubsection{单元一｜湿地：法律可以划界，水却会继续流动}\label{ux5355ux5143ux4e00ux6e7fux5730ux6cd5ux5f8bux53efux4ee5ux5212ux754cux6c34ux5374ux4f1aux7ee7ux7eedux6d41ux52a8-1}

\textbf{W1}　P1
的三个动词分别对应``储水、过滤、保护''。第36题A概括其价值，而不是把生态环境写成商业潜力。

\textbf{W2}　法院确立了更严格的连接条件，使EPA可监管的部分湿地范围缩小；这不等于各州法或全部其他保护消失。

\textbf{W3}　因为流域跨越多个州和地区，其他地方的保护选择会影响海湾；本州法律不足以控制上游全部行为。

\textbf{W4}　该例说明跨州外部影响使协调必要；原文没有断言所有农民主观恶意。

\textbf{W5}　作者面向州立法者，主张认真审视本州湿地保护，回应前文对跨境环境影响的担忧。

\subsubsection{单元二｜野火：少发生的火，为什么可能留下更多风险}\label{ux5355ux5143ux4e8cux91ceux706bux5c11ux53d1ux751fux7684ux706bux4e3aux4ec0ux4e48ux53efux80fdux7559ux4e0bux66f4ux591aux98ceux9669-1}

\textbf{F1}　他意外发现殖民前已有反复发生的火灾记录；任职职责只是叙事铺垫。

\textbf{F2}　作者并列多个生态效应来解释低烈度火灾的益处，而非仅陈述某项树木治疗。

\textbf{F3}　关注所研究的点位、火灾频率及过去几十年相对于历史记录的变化；不能泛化所有地区严重程度。

\textbf{F4}　成功扑灭局部火灾与改变长期燃料状况是不同时间尺度和不同结果；作者承认前者，担忧后者。

\textbf{F5}　并非如此；论证依赖适当环境、可控制条件与减轻后续风险的目的。

\textbf{F6}　共同支持在适当条件下利用计划烧除；Prichard突出了气候变化使该工作更加困难。

\subsection{三　历年真题的参考答案与选项比较}\label{ux4e09-ux5386ux5e74ux771fux9898ux7684ux53c2ux8003ux7b54ux6848ux4e0eux9009ux9879ux6bd4ux8f83}

以下是以试卷文字及独立转录交叉核对的\textbf{非官方参考答案}。涉及现实政策及历史研究时，只说明文章自己的论证；不把作者立场当作已证实的普遍规律。

\subsubsection{2024 年英语（一） Reading Text
4（36---40题）}\label{ux5e74ux82f1ux8bedux4e00-reading-text-43640ux9898}

\textbf{第36题　A　｜段内概括}

\textbf{为何选择}：P1列举储水、过滤与保护功能，最佳选项概括自然环境的价值。

\textbf{其余选项的偏差}：B称争议保护区、C强调商业潜力、D说是社区负担，均没有保留P1正面列举的生态功能。

\textbf{第37题　B　｜制度变化}

\textbf{为何选择}：P2指出联邦环境保护机构监管湿地的权限被收窄；答案不是所有法规取消。

\textbf{其余选项的偏差}：A说强化治理与方向相反；C说消除冲突、D说矿业经营者反对，不是本段证据。

\textbf{第38题　A　｜态度判断}

\textbf{为何选择}：P3先说州内有保护，再说跨州流域使这种安心短视；作者对未来感到担忧。

\textbf{其余选项的偏差}：B困惑、C放心、D鼓励，都与``其他州是否同样保护''的担忧不符。

\textbf{第39题　D　｜有限推断}

\textbf{为何选择}：P4强调联邦机构能超越相邻辖区局部利益，农业径流的例子说明跨地区协调的重要性。

\textbf{其余选项的偏差}：A把协调扩大成已经``恢复辖区平衡''；B扯到渔业改革；C将跨境协调改写成仅尊重州权，均改变了原论点。

\textbf{第40题　B　｜作者呼吁}

\textbf{为何选择}：P5明确呼吁州立法者重新审视湿地保护，准确概括为重视保护。

\textbf{其余选项的偏差}：A聚焦限制地主影响、C狭窄化为扩大避难所、D聚焦濒危物种福利，均不是末段主要行动。

\subsubsection{2026 年英语（一） Reading Text
4（36---40题）}\label{ux5e74ux82f1ux8bedux4e00-reading-text-43640ux9898-1}

\textbf{第36题　D　｜人物发现}

\textbf{为何选择}：P1惊讶点是前殖民时期野火反复发生，树轮火痕提供历史证据。

\textbf{其余选项的偏差}：A缺乏研究、B古代欧洲灭火、C1970年代管理实践均把背景换成惊讶对象。

\textbf{第37题　C　｜段落主旨}

\textbf{为何选择}：P2并列介绍低烈度火清理落叶、减少病害、促进更新等多项生态益处。

\textbf{其余选项的偏差}：A起因、B某一种树的治疗、D笼统生态重要性，不能概括该段机制。

\textbf{第38题　B　｜数据指标}

\textbf{为何选择}：P3的研究指出特定监测点近期火灾次数相对历史基线显著减少。

\textbf{其余选项的偏差}：A严重程度波动、C威胁被低估、D资料误读都更换了研究测量对象。

\textbf{第39题　D　｜机制推断}

\textbf{为何选择}：P4同时讨论长期扑火造成的燃料积累与气候变化，二者可能使极端条件下的火更难控制。

\textbf{其余选项的偏差}：A训练不足、B公众质疑均无证据；C说传统措施``无用''程度过强。

\textbf{第40题　C　｜多人共识}

\textbf{为何选择}：P5---6两位专家都支持有条件的计划用火；承认难度并不等于反对。

\textbf{其余选项的偏差}：A预测困难不是共同论点；B仅评估而未实施，改变倡议；D加倍投入探测也不是共同主张。

\subsection{四　句法工坊的结构与译文}\label{ux56db-ux53e5ux6cd5ux5de5ux574aux7684ux7ed3ux6784ux4e0eux8bd1ux6587}

\textbf{S1}　法律边界的变化不是自然地理边界的变化。

参考译文：当一项规则的边界发生变化时，河流本身仍然没有改变。

\textbf{S2}　两命题并非直接矛盾：本地行动不能必然控制外来输入。

参考译文：该州已经采取严格管控，但仍无法阻止污染从上游进入。

\textbf{S3}　at these sites / frequency是必要的范围约束。

参考译文：证据表明这些点位的火灾频率下降，而不是全球火灾危险都在下降。

\textbf{S4}　作者比较``数量''与``长期效果''，并非完全否定干预数量指标。

参考译文：指导政策的不应仅是干预次数，还应是其长期效果。

\textbf{S5}　是否安全是有待核查的问题，不是已确定事实。

参考译文：由于计划用火被推迟，团队有时间核查条件是否安全。

\textbf{S6}　让步不否定即时益处；对长期风险作可能性判断。

参考译文：即便一项干预减少了即时损失，它仍可能造成多年后才显现的风险。

\subsection{五　原创阅读：答案、证据与干扰分析}\label{ux4e94-ux539fux521bux9605ux8bfbux7b54ux6848ux8bc1ux636eux4e0eux5e72ux6270ux5206ux6790}

\subsubsection{Reading A · The Quiet Value of a
Floodplain}\label{reading-a-the-quiet-value-of-a-floodplain-1}

\textbf{A1　B}　P1明确保护在一地有效却将部分水流推向下游。A使用every夸大，C否定全文前提，D原文无据。

\textbf{A2　C}　P2强调单季收入依赖和暂时损失的压力；其他项没有对应证据。

\textbf{A3　A}　作者比较两种成本边界，结论随计入项目变化；不是时间事件。

\textbf{A4　D}　P4明确fund for affected
farms；不能推出完全消除洪灾或保证利润。

\textbf{A5　A}　文章不是反对堤墙而是强调跨境效应与公平分担；其余选项绝对化或夸大工具效果。

\textbf{AO1}　同时考虑城镇防洪收益、农场损失与补偿，并通过流域报告使影响可比较。

\textbf{AO2}　不能；它使成本公开、判断成为可能，但公平仍取决于实际分担。

\textbf{本文论证复盘}：先找P1问题提出，再看中段出现的相反考量，最后将末段的限定结论写成一句``不超过原文证据范围''的话。

\subsubsection{Reading B · The Fires a Park Did Not
Have}\label{reading-b-the-fires-a-park-did-not-have-1}

\textbf{B1　A}　P1指出单一评估指标为事件数量。

\textbf{B2　B}　温和与极端火灾作用不同，故不能只用发生次数评价。

\textbf{B3　C}　P3列出风/湿度/道路检查，限制性计划体现审慎。

\textbf{B4　D}　P3明确风险不能归零，其余过强或原文相反。

\textbf{B5　B}　P4将评价框架从单一次数转向燃料/损害/成本/试验影响。

\textbf{BO1}　不考虑火灾严重程度、长期可燃物积累和后续风险。

\textbf{BO2}　没有；重点是有条件、有监测的试验与多指标评估。

\textbf{本文论证复盘}：先找P1问题提出，再看中段出现的相反考量，最后将末段的限定结论写成一句``不超过原文证据范围''的话。

\subsection{六　专项练习答案}\label{ux516d-ux4e13ux9879ux7ec3ux4e60ux7b54ux6848}

\subsubsection{选词填空
V01---V08}\label{ux9009ux8bcdux586bux7a7a-v01v08}

\begin{itemize}
\item
  \textbf{V01}　narrow。注意用法是否符合该句语法，不能只凭中文近义判断。
\item
  \textbf{V02}　runoff。注意用法是否符合该句语法，不能只凭中文近义判断。
\item
  \textbf{V03}　jurisdictions。注意用法是否符合该句语法，不能只凭中文近义判断。
\item
  \textbf{V04}　compiled。注意用法是否符合该句语法，不能只凭中文近义判断。
\item
  \textbf{V05}　erode。注意用法是否符合该句语法，不能只凭中文近义判断。
\item
  \textbf{V06}　call
  on。注意用法是否符合该句语法，不能只凭中文近义判断。
\item
  \textbf{V07}　prescribed。注意用法是否符合该句语法，不能只凭中文近义判断。
\item
  \textbf{V08}　double-edged
  sword。注意用法是否符合该句语法，不能只凭中文近义判断。
\end{itemize}

\subsubsection{篇章衔接 G1---G2}\label{ux7bc7ux7ae0ux8854ux63a5-g1g2}

\begin{itemize}
\item
  \textbf{G1}　A。G1承接容量数据，引出``谁支付成本''的问题，使下文财政分割自然出现。
\item
  \textbf{G2}　C。G2由独立预算导出协调必要性，再衔接新协议的成本分担。
\end{itemize}

\subsubsection{英译汉 T01---T04}\label{ux82f1ux8bd1ux6c49-t01t04}

\textbf{T01}　虽然该县改善了本地防洪设施，污染仍从管辖范围外的地区不断进入。

\textbf{观察要点}　although让步；jurisdiction的抽象范围。

\textbf{T02}　这些证据涉及特定点位小型火灾的发生频率，而非所有野火的严重程度。

\textbf{观察要点}　concerns限定主题；rather than区分指标。

\textbf{T03}　一种降低即时危险的方法，如果改变了森林条件，仍可能带来长期风险。

\textbf{观察要点}　定语从句+may+if组合。

\textbf{T04}　与其笼统询问一项政策好坏，不如考察哪些群体受益、哪些群体承担成本。

\textbf{观察要点}　instead of+两个which并列。

\subsubsection{汉译英 P01---P02}\label{ux6c49ux8bd1ux82f1-p01p02}

\textbf{P01}　We need to specify the scope of an intervention rather
than extend a local success to an entire river basin.

\textbf{观察要点}　specify scope；rather than禁止过度泛化。

\textbf{P02}　Prescribed burning can reduce certain risks, but its use
depends on suitable weather conditions.

\textbf{观察要点}　can表示可能；certain限定；depend on表达条件。

\subsubsection{写作范文｜不要照抄，先判断它如何满足任务}\label{ux5199ux4f5cux8303ux6587ux4e0dux8981ux7167ux6284ux5148ux5224ux65adux5b83ux5982ux4f55ux6ee1ux8db3ux4efbux52a1}

\textbf{原创参考短文（约100词）}

An environmental policy should be judged by more than one number. A town
may report fewer floods after building a wall, yet villages downstream
may face greater risk. The local result is real, but it does not
describe the whole river basin. A fairer assessment would measure damage
across districts and include the costs borne by farmers and residents.
Reports should also compare short-term gains with long-term effects.
This does not mean that every flood barrier is misguided. It means that
success needs a clearly defined scope and several forms of evidence.
Without them, an improvement in one place may be mistaken for progress
everywhere.

\textbf{自己校改的三个问题}　是否区分具体数据与广泛结论？是否有连接词说明让步、转折与建议？不要求刻意使用生僻词，强调准确。

\textbf{复习建议}　首次订正时用已有线索重做错题；过几天只看英文新句，不看解析重述意思。间隔时间可随真实掌握情况调整，不能把``讲过''当``会了''。

\clearpage

\section{英中双语对照附录}\label{ux82f1ux4e2dux53ccux8bedux5bf9ux7167ux9644ux5f55}

本附录集中收录本册\textbf{自编英文}材料的中译，帮助做题后逐句核对。真题整篇原文及四选项请使用前述试卷链接查阅，本册不大量转录。

\subsection{A　入门诊断｜英中逐句}\label{a-ux5165ux95e8ux8bcaux65adux82f1ux4e2dux9010ux53e5}

\textbf{诊断 01}　A measure may protect local streams without reducing
pollution downstream.

\textbf{中译}　一项措施可能保护本地溪流，却未必减少下游污染。

\textbf{诊断 02}　The new rule narrowed the agency's authority; it did
not abolish the agency.

\textbf{中译}　新规缩小了该机构的权限，却没有撤销这一机构。

\textbf{诊断 03}　Although fewer small fires occurred, the forest
accumulated more dry fuel.

\textbf{中译}　虽然小规模火灾发生得更少，森林却积累了更多干燥可燃物。

\textbf{诊断 04}　Evidence from selected forests cannot establish a
universal trend.

\textbf{中译}　来自特定森林的证据不足以确立普遍趋势。

\subsection{B　单元一｜湿地：法律可以划界，水却会继续流动｜原创新句}\label{b-ux5355ux5143ux4e00ux6e7fux5730ux6cd5ux5f8bux53efux4ee5ux5212ux754cux6c34ux5374ux4f1aux7ee7ux7eedux6d41ux52a8ux539fux521bux65b0ux53e5}

\textbf{W1}　The marsh stores rainwater, filters waste and protects
nearby homes.

\textbf{中译}　沼泽蓄存雨水、过滤废弃物并保护附近的房屋。

\textbf{W2}　The court restricted the board's powers to a narrower group
of cases.

\textbf{中译}　法院将该委员会的权限限制在更窄的一类案件之内。

\textbf{W3}　A city may maintain strict standards, yet pollution can
arrive from outside it.

\textbf{中译}　一个城市可能维持严格标准，污染却仍可能从市界之外进入。

\textbf{W4}　A decision made upstream can impose costs on communities
that had no say in it.

\textbf{中译}　上游作出的决定可能让没有参与决策的下游社区承担成本。

\textbf{W5}　The editorial urges lawmakers to review safeguards before
the damage spreads.

\textbf{中译}　这篇社论敦促立法者在损害蔓延前重新审视保护措施。

\subsubsection{词汇例句的翻译}\label{ux8bcdux6c47ux4f8bux53e5ux7684ux7ffbux8bd1}

\textbf{filter}　Wetlands can filter some pollutants before water enters
the reservoir.

\textbf{中译}　湿地能在水流进入水库前过滤部分污染物。

\textbf{authority}　Local councils lack the authority to change national
safety standards.

\textbf{中译}　地方议会无权更改全国性的安全标准。

\textbf{scope}　The amendment reduced the scope of the inspection
programme.

\textbf{中译}　修正案缩小了检查计划的适用范围。

\textbf{at odds with}　The company's expansion plans were at odds with
the town's conservation goals.

\textbf{中译}　企业的扩张计划与城镇的保护目标相冲突。

\textbf{repercussion}　The abrupt change had repercussions for farms far
beyond the city.

\textbf{中译}　这一突然变化对城市之外很远的农场产生了连带影响。

\textbf{jurisdiction}　The river passes through several jurisdictions
with different rules.

\textbf{中译}　这条河流经数个实行不同规则的管辖区。

\textbf{shortsighted}　Cutting flood prevention funds may be
shortsighted even if it saves money today.

\textbf{中译}　削减防洪经费即使今天能省钱，也可能缺乏长远眼光。

\textbf{transcend}　Cooperation can transcend state borders when rivers
connect communities.

\textbf{中译}　当河流连接各地社区时，合作可以跨越州界。

\textbf{runoff}　Heavy rainfall carried fertilizer runoff into the
nearby lake.

\textbf{中译}　强降雨将含肥料的地表径流带入附近湖泊。

\textbf{call on}　Residents called on the council to publish its
water-quality data.

\textbf{中译}　居民呼吁议会公布水质数据。

\textbf{stakes}　The debate involves high stakes for fishing towns
downstream.

\textbf{中译}　这场争论关系到下游渔业城镇的重大利益。

\subsection{B　单元二｜野火：少发生的火，为什么可能留下更多风险｜原创新句}\label{b-ux5355ux5143ux4e8cux91ceux706bux5c11ux53d1ux751fux7684ux706bux4e3aux4ec0ux4e48ux53efux80fdux7559ux4e0bux66f4ux591aux98ceux9669ux539fux521bux65b0ux53e5}

\textbf{F1}　Old trees preserve evidence that a recent survey may
overlook.

\textbf{中译}　古树保留着近期调查可能忽略的证据。

\textbf{F2}　A mild burn can remove dry litter without destroying mature
trees.

\textbf{中译}　温和燃烧能够清除干燥落叶而不毁坏成树。

\textbf{F3}　At the sampled sites, recent fire counts were lower than
the historical pattern suggested.

\textbf{中译}　在抽样地点，近期火灾次数低于历史变化模式所显示的水平。

\textbf{F4}　Suppressing every small fire may allow dry fuel to
accumulate over time.

\textbf{中译}　扑灭所有小火可能让干燥可燃物随着时间积累。

\textbf{F5}　A prescribed burn is planned for conditions in which the
flames can be managed.

\textbf{中译}　计划烧除会安排在火势能够得到控制的条件下实施。

\textbf{F6}　The policy may become more necessary just as it becomes
harder to carry out safely.

\textbf{中译}　一项政策可能变得更必要，同时也变得更难安全实施。

\subsubsection{词汇例句的翻译}\label{ux8bcdux6c47ux4f8bux53e5ux7684ux7ffbux8bd1-1}

\textbf{mandate}　Her mandate was to reduce immediate threats to nearby
villages.

\textbf{中译}　她的职责是减少附近村庄面临的直接威胁。

\textbf{curb}　New measures have curbed the spread of small grass fires.

\textbf{中译}　新措施抑制了小规模草地火灾的蔓延。

\textbf{peter out}　The small fire petered out as the weather became
cooler.

\textbf{中译}　天气转凉后，小火逐渐熄灭了。

\textbf{compile}　Scientists compiled records from trees across the
region.

\textbf{中译}　科学家汇集了这一地区多棵树的记录。

\textbf{shortfall}　The survey found a shortfall in routine maintenance
compared with the previous decade.

\textbf{中译}　调查发现，与此前十年相比，常规维护出现了缺口。

\textbf{suppression}　Decades of suppression altered the conditions
under the forest canopy.

\textbf{中译}　数十年的灭火干预改变了林冠下的环境条件。

\textbf{testament}　The restored trail is a testament to the volunteers'
patience.

\textbf{中译}　修复后的步道体现了志愿者的耐心。

\textbf{erode}　Repeated delays eroded residents' trust in the warning
system.

\textbf{中译}　反复延误逐渐削弱了居民对预警系统的信任。

\textbf{ignite}　A single spark can ignite a fire in unusually dry
grass.

\textbf{中译}　在异常干燥的草地上，一粒火星就可能引发火灾。

\textbf{prescribed}　The team postponed the prescribed burn because the
wind shifted.

\textbf{中译}　由于风向变化，团队推迟了计划烧除。

\textbf{tactic}　Managers changed their tactic when the weather became
unpredictable.

\textbf{中译}　天气变得难以预测时，管理人员改变了具体手段。

\textbf{double-edged sword}　Improved access is a double-edged sword: it
supports tourism but increases pressure on habitats.

\textbf{中译}　改善交通是一把双刃剑：它支持旅游，却增加了栖息地压力。

\subsection{C　句法工坊｜六句双语}\label{c-ux53e5ux6cd5ux5de5ux574aux516dux53e5ux53ccux8bed}

\textbf{S1}　When the boundary of a rule changes, the river itself
remains unchanged.

\textbf{中译}　当一项规则的边界发生变化时，河流本身仍然没有改变。

\textbf{S2}　The state had adopted strict controls, yet it could not
prevent pollution arriving from upstream.

\textbf{中译}　该州已经采取严格管控，但仍无法阻止污染从上游进入。

\textbf{S3}　The evidence suggests a decline in fire frequency at these
sites, not a universal decline in fire danger.

\textbf{中译}　证据表明这些点位的火灾频率下降，而不是全球火灾危险都在下降。

\textbf{S4}　It is not the number of interventions alone but their
long-term effects that should guide policy.

\textbf{中译}　指导政策的不应仅是干预次数，还应是其长期效果。

\textbf{S5}　Because the planned fire was delayed, the team had time to
check whether conditions were safe.

\textbf{中译}　由于计划用火被推迟，团队有时间核查条件是否安全。

\textbf{S6}　Even if an intervention reduces immediate losses, it may
create risks that appear only years later.

\textbf{中译}　即便一项干预减少了即时损失，它仍可能造成多年后才显现的风险。

\subsection{D1　Reading A · The Quiet Value of a
Floodplain｜全文双语}\label{d1-reading-a-the-quiet-value-of-a-floodplainux5168ux6587ux53ccux8bed}

\textbf{P1（英文）}　A town on the River Lorne spent years strengthening
the wall that separated its shopping district from the water. The wall
worked during ordinary storms. Shops stayed open, and local officials
described the project as a success. But after an unusually wet winter,
engineers discovered that the wall had pushed more water towards
villages downstream. Protection in one place had changed the
distribution of danger, rather than removed it.

\textbf{P1（中译）}　洛恩河畔的一座城镇花了数年加固商业区与河流之间的堤墙。普通风暴中，这道墙发挥了作用，商店得以照常营业，地方官员称项目成功。但在一个异常多雨的冬季之后，工程师发现堤墙把更多河水推向了下游村庄。一地的保护改变了风险的分布，却没有消除风险。

\textbf{P2（英文）}　The regional council then considered restoring a
floodplain north of the town. Under the proposal, part of the farmland
would be allowed to flood during periods of heavy rain. Supporters
argued that temporary flooding on open land could prevent far greater
damage in densely populated areas. Farmers, however, pointed out that
even temporary losses could place serious pressure on families whose
income depended on a single season.

\textbf{P2（中译）}　地区议会随后考虑恢复城镇北侧的一片洪泛区。按该方案，暴雨期间允许部分农田暂时被淹。支持者认为，开放土地上的暂时积水可以避免人口稠密地区遭受更大的损失。但农民指出，即使是暂时损失，也可能对依靠单个种植季收入的家庭造成沉重压力。

\textbf{P3（英文）}　Neither side disputed the need for flood
protection. Their disagreement concerned who should bear its costs and
how those costs should be measured. If the council counted only repairs
to buildings, the floodplain looked inexpensive. If it included lost
crops, disrupted transport and the uncertainty faced by farmers, the
calculation became more complicated.

\textbf{P3（中译）}　双方都不否认防洪的必要性。他们争论的是谁该承担成本、成本应如何计算。如果议会只计算建筑修缮费用，洪泛区方案似乎很便宜；如果把作物损失、交通中断和农民面对的不确定性也纳入考虑，计算就变得复杂得多。

\textbf{P4（英文）}　The final plan combined smaller walls with a
compensation fund for affected farms. It also required the council to
publish annual reports on flooding across the whole river basin. The
arrangement did not guarantee that every household would be safe. It
did, however, make the trade-offs visible and allow communities to judge
whether protection was being shared fairly.

\textbf{P4（中译）}　最终方案将较小的堤墙与对受影响农场的补偿基金结合起来，并要求议会每年公布整个流域的洪灾报告。该安排无法保证每户家庭都绝对安全，但确实让取舍变得可见，使各社区能够判断保护是否得到公平分担。

\subsubsection{阅读题及全部选项英中对照}\label{ux9605ux8bfbux9898ux53caux5168ux90e8ux9009ux9879ux82f1ux4e2dux5bf9ux7167}

\textbf{A1.}　What did the engineers discover after the wet winter?

\textbf{中译}　那个多雨冬季之后，工程师发现了什么？

\begin{itemize}
\item
  \textbf{A.} The wall had failed in every storm.

  \textbf{中译}　堤墙在每场风暴中都失效。
\item
  \textbf{B.} The wall had shifted some risk downstream.

  \textbf{中译}　堤墙将部分风险转移到了下游。
\item
  \textbf{C.} The town had no need for flood protection.

  \textbf{中译}　城镇已不需要防洪。
\item
  \textbf{D.} The villages had built a second wall.

  \textbf{中译}　村庄修建了第二道堤墙。
\end{itemize}

\textbf{A2.}　The farmers were mainly concerned about

\textbf{中译}　农民主要担忧什么？

\begin{itemize}
\item
  \textbf{A.} the cost of replacing the town's shops.

  \textbf{中译}　更换城镇商铺的费用。
\item
  \textbf{B.} the absence of an annual report.

  \textbf{中译}　缺少年报。
\item
  \textbf{C.} the pressure of seasonal losses on family income.

  \textbf{中译}　季节性损失对家庭收入的压力。
\item
  \textbf{D.} the closing of the regional council.

  \textbf{中译}　地区议会关闭。
\end{itemize}

\textbf{A3.}　The third paragraph mainly examines

\textbf{中译}　第三段主要分析什么？

\begin{itemize}
\item
  \textbf{A.} how costs are defined and counted.

  \textbf{中译}　成本如何定义和计算。
\item
  \textbf{B.} why farmers oppose all public spending.

  \textbf{中译}　农民为何反对所有公共支出。
\item
  \textbf{C.} whether shop repairs are legally required.

  \textbf{中译}　商店修缮是否具有法律强制性。
\item
  \textbf{D.} when the river changed its course.

  \textbf{中译}　河流何时改道。
\end{itemize}

\textbf{A4.}　The compensation fund was intended to

\textbf{中译}　补偿基金旨在什么？

\begin{itemize}
\item
  \textbf{A.} eliminate the possibility of flooding.

  \textbf{中译}　排除任何洪水可能。
\item
  \textbf{B.} replace all flood reports.

  \textbf{中译}　取代所有洪灾报告。
\item
  \textbf{C.} guarantee profits to every business.

  \textbf{中译}　确保所有企业盈利。
\item
  \textbf{D.} address losses borne by affected farms.

  \textbf{中译}　回应受影响农场承担的损失。
\end{itemize}

\textbf{A5.}　Which statement best captures the author's point?

\textbf{中译}　哪项最能概括作者要点？

\begin{itemize}
\item
  \textbf{A.} Good protection requires attention to where costs fall.

  \textbf{中译}　良好的保护需要关注成本落在谁身上。
\item
  \textbf{B.} Local walls should never be built.

  \textbf{中译}　地方堤墙绝不应该建。
\item
  \textbf{C.} Farmers should determine all public policy.

  \textbf{中译}　农民应决定一切公共政策。
\item
  \textbf{D.} Annual reports can prevent natural disasters.

  \textbf{中译}　年报可以防止自然灾害。
\end{itemize}

\subsubsection{无选项简答的题干对照}\label{ux65e0ux9009ux9879ux7b80ux7b54ux7684ux9898ux5e72ux5bf9ux7167}

\textbf{O1}　Why does the author call the final plan a trade-off?

\textbf{中译}　为什么作者把最终方案看作一种取舍？

\textbf{O2}　Does publication of reports itself ensure fairness?

\textbf{中译}　仅仅公布报告就能确保公平吗？

\subsection{D2　Reading B · The Fires a Park Did Not
Have｜全文双语}\label{d2-reading-b-the-fires-a-park-did-not-haveux5168ux6587ux53ccux8bed}

\textbf{P1（英文）}　For decades, staff at Alder Park judged their fire
programme by a simple number: how few fires occurred each summer. The
fewer incidents they recorded, the more successful the programme
appeared. Visitors felt reassured, and managers could report that fewer
people had been exposed to smoke. Yet the measure offered no information
about what was happening to the forest floor.

\textbf{P1（中译）}　多年来，奥尔德公园的工作人员用一个简单数字评价防火计划：每个夏季火灾发生得有多少次。记录越少，计划看起来就越成功。游客感到安心，管理者也能报告接触烟雾的人减少了。然而，这个指标没有提供森林地面状况的信息。

\textbf{P2（英文）}　A survey later found thick layers of dry branches
in areas where small burns had long been suppressed. Ecologists did not
argue that every wildfire was beneficial. They distinguished mild fires,
which might clear some fallen material, from extreme fires capable of
damaging mature trees and nearby buildings. Treating both as a single
category had hidden a difference important to future decisions.

\textbf{P2（中译）}　后来的调查发现，长期压制小型燃烧的区域积累了厚厚的干树枝。生态学家并未声称所有野火都有益。他们区分可能清理部分落枝的温和火灾与足以损坏成树和附近建筑的极端火灾。把两者归为同一类，遮蔽了一个对未来决策重要的差异。

\textbf{P3（英文）}　The park proposed a limited trial of planned
burning. Fires would be introduced only on selected days, after checks
of wind, moisture and access roads. A small team would monitor the
operation, and nearby residents would be informed in advance. Critics
reasonably asked whether changing weather could turn a controlled burn
into an emergency. The project team acknowledged that this risk could
not be reduced to zero.

\textbf{P3（中译）}　公园提出进行有限的计划烧除试验。只在选定日期、检查风、湿度和道路条件后才点火；一个小组负责监测，附近居民将提前获知。批评者合理地询问：天气变化是否可能将可控燃烧变成紧急事件？项目组承认，这一风险不能降为零。

\textbf{P4（英文）}　For this reason, the trial would be evaluated by
several measures rather than a single fire count: fuel levels, damage
from later fires, costs of preparation and any harm caused by the trial
itself. The aim was not to celebrate fire. It was to learn which
combination of prevention and carefully managed burning could make the
park safer over time.

\textbf{P4（中译）}　因此，试验将采用多种指标而非单一火灾次数评价，包括可燃物水平、之后火灾的损害、准备成本和试验自身造成的任何损失。目的不是赞美火，而是弄清防火与谨慎用火怎样组合，才能使公园长期更加安全。

\subsubsection{阅读题及全部选项英中对照}\label{ux9605ux8bfbux9898ux53caux5168ux90e8ux9009ux9879ux82f1ux4e2dux5bf9ux7167-1}

\textbf{B1.}　The park originally assessed success mainly through

\textbf{中译}　公园最初主要依据什么评价成功？

\begin{itemize}
\item
  \textbf{A.} the number of summer fires.

  \textbf{中译}　夏季火灾次数。
\item
  \textbf{B.} the age of its oldest trees.

  \textbf{中译}　最古老树木的年龄。
\item
  \textbf{C.} the amount of rain in winter.

  \textbf{中译}　冬季降雨量。
\item
  \textbf{D.} the cost of public transport.

  \textbf{中译}　公共交通成本。
\end{itemize}

\textbf{B2.}　The ecologists' distinction in Paragraph 2 suggests that

\textbf{中译}　第二段生态学家的区分说明什么？

\begin{itemize}
\item
  \textbf{A.} all fires should burn freely.

  \textbf{中译}　所有火灾都应自由燃烧。
\item
  \textbf{B.} fire severity affects how an event is assessed.

  \textbf{中译}　火灾严重程度影响事件评价。
\item
  \textbf{C.} smoke is always harmless.

  \textbf{中译}　烟雾总是无害。
\item
  \textbf{D.} old trees prevent extreme fires.

  \textbf{中译}　老树可防止极端火灾。
\end{itemize}

\textbf{B3.}　Why would the trial be limited to selected days?

\textbf{中译}　为什么试验只在选定日期进行？

\begin{itemize}
\item
  \textbf{A.} To maximise the total number of fires.

  \textbf{中译}　最大化火灾总次数。
\item
  \textbf{B.} To avoid telling residents.

  \textbf{中译}　为了不告知居民。
\item
  \textbf{C.} To test risks under assessed conditions.

  \textbf{中译}　在经评估的条件下控制风险。
\item
  \textbf{D.} To stop all forest surveys.

  \textbf{中译}　停止所有森林调查。
\end{itemize}

\textbf{B4.}　The project team conceded that planned burning

\textbf{中译}　项目组承认计划烧除什么局限？

\begin{itemize}
\item
  \textbf{A.} had already destroyed every building.

  \textbf{中译}　已摧毁所有建筑。
\item
  \textbf{B.} would never be useful.

  \textbf{中译}　永远无效。
\item
  \textbf{C.} required no monitoring.

  \textbf{中译}　无需监测。
\item
  \textbf{D.} could not be made completely risk-free.

  \textbf{中译}　不能完全消除风险。
\end{itemize}

\textbf{B5.}　The central argument is that managers should

\textbf{中译}　全文中心观点是管理者应当如何？

\begin{itemize}
\item
  \textbf{A.} replace all prevention with fires.

  \textbf{中译}　用火取代所有预防。
\item
  \textbf{B.} judge measures by several long-term outcomes.

  \textbf{中译}　综合若干长期结果评价措施。
\item
  \textbf{C.} avoid listening to local residents.

  \textbf{中译}　避免听取居民意见。
\item
  \textbf{D.} prefer immediate numbers to consequences.

  \textbf{中译}　优先看即时数字而非后果。
\end{itemize}

\subsubsection{无选项简答的题干对照}\label{ux65e0ux9009ux9879ux7b80ux7b54ux7684ux9898ux5e72ux5bf9ux7167-1}

\textbf{O1}　What is the weakness of counting fires alone?

\textbf{中译}　只数火灾次数有什么局限？

\textbf{O2}　Does the author recommend uncontrolled wildfires?

\textbf{中译}　作者是否建议不加控制的野火？

\subsection{E　专项练习英中对照}\label{e-ux4e13ux9879ux7ec3ux4e60ux82f1ux4e2dux5bf9ux7167}

\subsubsection{完成后的词汇句子}\label{ux5b8cux6210ux540eux7684ux8bcdux6c47ux53e5ux5b50}

\textbf{V01}　The proposal would narrow the agency's powers to a smaller
set of cases.

\textbf{中译}　这项提案将把该机构的权限缩小到更少的一类案件。

\textbf{V02}　Rainwater carried agricultural runoff into the river after
the storm.

\textbf{中译}　暴风雨后，雨水将农业地表径流带入河流。

\textbf{V03}　The river crosses several jurisdictions, each with its own
rules.

\textbf{中译}　这条河流经数个辖区，每个辖区都有自己的规则。

\textbf{V04}　Scientists compiled fire records collected from hundreds
of trees.

\textbf{中译}　科学家汇编了从数百棵树采集的火灾记录。

\textbf{V05}　A series of small cuts may gradually erode public trust.

\textbf{中译}　一系列小幅削减可能逐渐损害公众信任。

\textbf{V06}　The council decided to call on residents to take part in
the review.

\textbf{中译}　议会决定呼吁居民参加审查。

\textbf{V07}　The carefully prescribed fire was delayed because the wind
changed.

\textbf{中译}　因为风向变化，精心安排的计划用火被推迟。

\textbf{V08}　A rule can be a double-edged sword when it protects one
group but burdens another.

\textbf{中译}　一项规则在保护某个群体却增加另一群体负担时，可能是一把双刃剑。

\subsubsection{篇章衔接的原句、备选句及中译}\label{ux7bc7ux7ae0ux8854ux63a5ux7684ux539fux53e5ux5907ux9009ux53e5ux53caux4e2dux8bd1}

\textbf{英文}　The committee first measured how much rainwater each
district could store. It reported that the largest upstream wetland
could hold a significant amount during heavy storms.

\textbf{中译}　委员会首先测量各区能够蓄存多少雨水，报告指出，最大的上游湿地能够在暴雨期间蓄存可观水量。

\textbf{英文}　Officials then discovered that local budgets were
prepared separately, even though flooding crossed district boundaries.

\textbf{中译}　官员随后发现，虽然洪水跨越各区边界，地方预算却是彼此独立编制的。

\textbf{英文}　The new agreement therefore shared monitoring costs and
required annual public reports.

\textbf{中译}　因此，新协议分担监测费用，并要求每年向公众报告。

\textbf{A}　Yet this storage capacity did not explain who should pay for
its protection.

\textbf{中译}　然而，这种蓄水能力并未说明谁该为保护它付费。

\textbf{B}　The scientists refused to measure any rainfall at all.

\textbf{中译}　科学家拒绝测量任何降雨。

\textbf{C}　Without coordination, no district could be sure how its
neighbours would act.

\textbf{中译}　如果缺乏协调，任何一个区都无法确定邻区会怎样行动。

\textbf{D}　The upstream wetland was built by the committee only
yesterday.

\textbf{中译}　这片上游湿地直到昨天才由委员会建成。

\subsubsection{翻译、表达与写作任务}\label{ux7ffbux8bd1ux8868ux8fbeux4e0eux5199ux4f5cux4efbux52a1}

\textbf{T01}　Although the county improved its own flood barriers,
pollution continued to arrive from areas outside its jurisdiction.

\textbf{中译}　虽然该县改善了本地防洪设施，污染仍从管辖范围外的地区不断进入。

\textbf{T02}　The evidence concerns the frequency of small fires at
selected sites rather than the severity of every wildfire.

\textbf{中译}　这些证据涉及特定点位小型火灾的发生频率，而非所有野火的严重程度。

\textbf{T03}　A method that reduces immediate danger may nevertheless
create a long-term risk if it changes the conditions of the forest.

\textbf{中译}　一种降低即时危险的方法，如果改变了森林条件，仍可能带来长期风险。

\textbf{T04}　Instead of asking whether a policy is good or bad in
general, we should examine which groups benefit and which bear the
costs.

\textbf{中译}　与其笼统询问一项政策好坏，不如考察哪些群体受益、哪些群体承担成本。

\textbf{P01}　We need to specify the scope of an intervention rather
than extend a local success to an entire river basin.

\textbf{中译}　必须说明干预措施的适用范围，不能把局部成功推广到全部流域。

\textbf{P02}　Prescribed burning can reduce certain risks, but its use
depends on suitable weather conditions.

\textbf{中译}　计划用火有助于减少某些风险，但其实施必须依赖适当的天气条件。

\textbf{写作任务（英文）}　Write about 90--120 words for a school
newsletter. Explain why one environmental indicator may be insufficient
for judging the success of a policy. Use an example and suggest a fairer
way of assessment.

\textbf{中译}　为校园通讯写约90---120词短文，说明为什么单一环境指标不足以判断政策成败，给出一个例子并提出更合理的评价方式。

\textbf{原创写作参考（英文）}　An environmental policy should be judged
by more than one number. A town may report fewer floods after building a
wall, yet villages downstream may face greater risk. The local result is
real, but it does not describe the whole river basin. A fairer
assessment would measure damage across districts and include the costs
borne by farmers and residents. Reports should also compare short-term
gains with long-term effects. This does not mean that every flood
barrier is misguided. It means that success needs a clearly defined
scope and several forms of evidence. Without them, an improvement in one
place may be mistaken for progress everywhere.

\textbf{中译}　环境政策不应只用一个数字来评价。一座城镇修建堤墙后可能报告洪灾更少，但下游村庄的风险却可能升高。本地结果是真实的，却无法描述整个流域。更公平的评估应测量不同辖区的损害，并纳入农民和居民承担的成本。报告还应比较短期收益和长期影响。这并不是说每道防洪堤都错误，而是说成功必须有清楚的适用范围和多种证据；否则，一地的改善可能被误认为处处都取得了进步。

\clearpage

\section{学习材料与出处}\label{ux5b66ux4e60ux6750ux6599ux4e0eux51faux5904}

本册真实试卷伴读依据相应年份英语（一）真题的公开整理版，并结合独立转录核对段落与第26---30或36---40题。参考答案仅为核对后的非官方参考答案。不同网站出现拼写、删节或翻译差异时，以试卷版对照，不把来源网页中的解析直接当成标准答案。

\begin{itemize}
\item
  2024 年英语（一） Reading Text
  4（36---40题）：\href{https://ac-wiki.org/general-skills/study/assets/2024-201.pdf}{打开试卷页面}
\item
  2026 年英语（一） Reading Text
  4（36---40题）：\href{https://english-exam.lazynote.cn/kaoyan/sections/2026-english-one/section2-part-a-4/}{打开试卷页面}
\end{itemize}

课外词汇学习参考：I. S. P. Nation, \emph{Learning Vocabulary in Another
Language}, Cambridge University
Press（2022），词汇知识不仅包括释义，还包括形式、搭配、语法功能与使用限制。

\textbf{使用提醒}：自编短文只用于训练，不冒充命题者写作或某一年试卷原文。它们的题目故意包含不同类型的错误项，以训练根据具体证据而不是根据口诀判断。

\end{document}
